% !TEX program = pdflatex
\documentclass[12pt]{article}
\usepackage[a4paper,margin=1in]{geometry}
\usepackage[T1]{fontenc}
\usepackage[utf8]{inputenc}
\usepackage{lmodern}
\usepackage{amsmath,amssymb,amsthm,mathtools,bm}
\usepackage{booktabs,array,tabularx}
\usepackage{enumitem}
\newcolumntype{L}[1]{>{\raggedright\arraybackslash}p{#1}}
\newcolumntype{Y}{>{\raggedright\arraybackslash}X}
\usepackage{setspace}
\usepackage[round,authoryear]{natbib}
\usepackage[colorlinks=true,linkcolor=blue,citecolor=blue,urlcolor=blue]{hyperref}

\onehalfspacing
\setlength{\parindent}{1.5em}
\setlength{\parskip}{0.12em}
\setlist[itemize]{leftmargin=2em,itemsep=0.25em,topsep=0.35em}
\setlist[enumerate]{leftmargin=2em,itemsep=0.25em,topsep=0.35em}
\emergencystretch=3em

\newtheorem{proposition}{Proposition}
\newtheorem{corollary}{Corollary}
\newtheorem{remark}{Remark}
\newtheorem{definition}{Definition}
\newtheorem{assumption}{Assumption}

\newcommand{\E}{\mathbb{E}}
\newcommand{\1}{\mathbf{1}}
\newcommand{\dd}{\mathrm{d}}
\newcommand{\calR}{\mathcal{R}}

\title{Regulatory Separation and Pharmaceutical Innovation: Evidence from China's MAH Reform}
\author{}
\date{}

\begin{document}
\maketitle
\thispagestyle{plain}

\begin{abstract}
\noindent
This paper develops a partial-equilibrium mechanism model of how China's
Marketing Authorization Holder (MAH) system changes the value of advancing
viable drug projects and their commercialization organization. Developers
choose one common original-drug innovation investment / project-advancement
intensity before project characteristics and organizational routes are
realized. MAH makes retained entrusted manufacturing legally available while
the developer retains authorization-holder responsibility. Developers differ
in project-advancement capability and internal manufacturing capability.
Internal and entrusted production use distinct technologies and organizational
costs. The new route is especially valuable when weak internal manufacturing
capability makes vertical integration costly. Route choice remains
deterministic, and the price of qualified manufacturing services is
endogenous. Reform effects are confined to project--developer pairs for which
the new route improves on the best old option; service scarcity can attenuate
these gains. Financing frictions may amplify this organizational mechanism,
but financing is not a baseline policy channel. The baseline does not treat
MAH as a demand, credit-supply, scientific-productivity, downstream-success,
or patenting shock.
\end{abstract}

\noindent\textit{Keywords:} MAH; pharmaceutical innovation; commercialization frictions; complementary assets; entrusted production; organizational choice.\\
\noindent\textit{JEL codes:} D22; L65; O31; O38; I18.

\section{Introduction}

China's MAH system separates the right to hold marketing authorization from a
strict requirement of in-house production. This institutional change is
potentially important for drug development because many research-originating
firms can generate valuable projects but lack complete manufacturing
capacity. The central mechanism studied here is therefore not demand expansion
or scientific productivity. It is an expected-commercialization-value
channel: when a firm anticipates that an advancing project can later use
qualified external manufacturing while the firm retains authorization, the
expected value of advancing that project may increase. The reform can be
particularly important when weak internal manufacturing capability makes
vertical integration costly. A developer without suitable existing capacity
can in principle build, acquire, validate, or adapt compliant production, but
that internalization commitment can be large and front-loaded. The baseline
captures this burden through route-specific setup and operating costs rather
than modeling financial intermediation. A financing-friction extension is
considered separately.

The mechanism builds on two established ideas. First, pharmaceutical
innovation responds to expected downstream rewards
\citep{Schmookler1966,AcemogluLinn2004,Finkelstein2004,DuboisDeMouzonScottMortonSeabright2015,BudishRoinWilliams2015}.
Second, an innovator's ability to appropriate returns depends on complementary
assets, firm boundaries, and access to specialized inputs
\citep{Teece1986,AroraFosfuriGambardella2001,GansStern2003,GrossmanHelpman2002,BolerMoxnesUlltveitMoe2015,Ma2026}.
MAH combines these ideas narrowly: it changes how an innovator can access
manufacturing capability while retaining the authorization asset. It does not
shift credit supply.

Recent evidence on China's pharmaceutical reforms helps separate this route
mechanism from two nearby channels. \citet{JiaMaYangZhang2023} study the 2015
drug-review reform and emphasize approval delay, review capacity, and the
composition of firms and products entering the pipeline.
\citet{BarwickXiaXia2026} study the National Reimbursement Drug List and
emphasize demand and market-size incentives. Those channels can coexist with
MAH, but they are not the MAH mechanism in this paper. The baseline holds
scientific productivity and the market-return environment fixed and asks
whether a regulated holder--producer separation route changes expected
commercialization value.

Closely related evidence also shows why project advancement and upstream
patenting must remain distinct. \citet{Gu2024} documents heterogeneous trial
and patent responses by financial resources. Those findings motivate a
financing-friction extension and empirical heterogeneity tests. The baseline
here isolates a different and more primitive MAH margin: the availability and
value of retained entrusted manufacturing relative to internal production,
transfer, and delay. This is a maintained model boundary, not an overstatement
of novelty; the baseline neither generates patents nor treats financing as the
policy channel.

For a development-economics audience, the relevant bottleneck is a regulated
complementary asset. MAH can relax a manufacturing bottleneck for
research-oriented firms while preserving holder responsibility. Scarcity in
qualified production support can still attenuate the effect.

The model is deliberately narrow. It is not a full dynamic structural model
of the pharmaceutical industry. It does not solve product-market clearing,
an input market for project advancement, free entry, exit, or an invariant
firm distribution. The baseline closes only the market directly tied to the
MAH mechanism: qualified manufacturing services. Reform-induced demand for
retained entrusted production can raise the equilibrium service price
\(p_m^*\), so the model does not mechanically imply a positive effect for
every project or developer.

The paper identifies a commercialization-organization channel of
pharmaceutical innovation. By separating marketing-authorization ownership
from physical manufacturing, MAH allows developers with weak internal
manufacturing capability to retain commercialization rights while purchasing
qualified production services. The model characterizes capability-based
organizational sorting, the project--developer pairs for which the new route
creates value, and the attenuation generated by scarcity in qualified
manufacturing capacity. Financing frictions may amplify the mechanism but are
not required for the baseline result.

The main text gives the shortest complete commercialization baseline. The
technical appendix records primitives, payoff accounting, pricing
derivations, deterministic sorting, project advancement, qualified-capacity
equilibrium, comparative statics, observed outcomes, and explicitly separated
inactive extensions, including commercialization financing.

\section{Institutional Background}

The relevant institutional point is timing. Under China's drug-registration
framework, a firm need not already be the final commercial holder when it
begins to advance a viable project. The economic choice is forward-looking: a
firm invests today while anticipating the manufacturing and commercialization
environment it will face if a project advances. Production support and
manufacturing readiness must be planned before final approval is observed.
After obtaining a Drug Approval License, the applicant becomes the MAH
\citep{NMPARegistration2022}. The model therefore treats MAH as changing the
anticipated availability and value of a future organizational route, not
current scientific productivity or downstream success.

The MAH reform creates a regulated holder-producer separation route. The 2016
pilot allowed eligible applicants in selected regions to become holders under
the pilot arrangement
\citep{StateCouncilMAHPilot2016,CFDAMAHImplementation2016}. The later legal
framework consolidated the principle that an MAH may manufacture by itself or
entrust production to a qualified manufacturer
\citep{DrugAdministrationLaw2019}. Subsequent NMPA rules require holder-side
quality responsibility, assessment of the entrusted manufacturer, quality and
entrustment agreements, and ongoing supervision
\citep{NMPAHolderResponsibility2022,NMPAManufacturing2022,NMPAEntrustedProduction2023}.

These rules imply two modeling restrictions. First, entrusted production is
not anonymous outsourcing or transfer. It is a regulated retained route in
which authorization remains with the developer while production is performed
by a qualified external manufacturer. Second, MAH does not eliminate holder
responsibility. The entrusted payoff therefore retains a holder-side burden
\(\mu_E\), while the institutional wedge \(\tau_E(M)\) governs effective
route availability. The downstream realization probability \(s(q)\) is
route independent and invariant to MAH in the baseline.

\section{Model}
\label{sec:p16m-model}

\subsection{Environment and timing}
\label{subsec:p16m-environment}

There is a continuum of developers. Developer \(i\) has predetermined
project-advancement capability \(a_i>0\), internal manufacturing capability
\(k_i>0\), with
\begin{equation}
  \theta_i=(a_i,k_i)\sim H(a,k).
  \label{eq:p16m-developer-type}
\end{equation}
The joint distribution permits arbitrary correlation. A viable project that
reaches commercialization-relevant planning draws value \(q>0\) and
manufacturing requirement \(m>0\) from the exogenous distribution \(F(q,m)\).
For empirical decompositions only,
\[
  F(q,m)=\sum_{g\in\{O,\mathrm{Inc}\}}\rho_gF_g(q,m),
  \qquad
  \rho_g\geq0,\quad \sum_g\rho_g=1.
\]
The class label \(g\) creates no class-specific control.

The institutional regime is \(M\in\{0,1\}\). The pre-MAH regime sets the
barrier to retained entrusted manufacturing at infinity; the MAH regime
makes it finite:
\begin{equation}
  \tau_E(0)=+\infty,
  \qquad
  \tau_E(1)=\bar\tau_E<+\infty,
  \qquad \bar\tau_E\geq0.
  \label{eq:p16m-policy}
\end{equation}
The available routes are internal manufacturing \(I\), retained entrusted
manufacturing \(E\), non-retained transfer \(T\), and abandonment or delay
\(A\). Under \(E\), the developer remains the authorization holder and a
qualified external producer manufactures. Thus \(E\) is not transfer.

Before project characteristics and routes are realized, developer \(i\)
chooses one common \(x_i\geq0\), interpreted as original-drug innovation
investment / project-advancement intensity. It is broader than pure clinical
development, but excludes basic research, compound discovery,
patent-generating effort, and patent applications. The expected measure of
viable projects reaching route planning is
\begin{equation}
  \lambda_i^{\mathrm{plan}}=a_ix_i.
  \label{eq:p16m-planning}
\end{equation}

Timing is as follows. First, \(M\) is observed and developers anticipate the
market-consistent qualified-service price \(p_m^*\). Second, after observing
\((a_i,k_i,M)\), they choose \(x_i\). Third, planning-stage projects draw
\((q,m)\) and choose among \(I,E,T,A\). Fourth, downstream realization occurs
with exogenous probability \(s(q)\). Finally, aggregate capacity demand and
supply must validate the anticipated service price. The last step is an
equilibrium consistency condition, not a late unanticipated shock. Upstream
research, patent generation, and financial intermediation are outside the
baseline decision problem.

\subsection{Commercialization technologies}
\label{subsec:p16m-technologies}

Conditional on successful commercialization, residual demand is
\(y(p;q)=Aqp^{-\varepsilon}\), where \(A>0\) and \(\varepsilon>1\).
Product-price optimization gives
\[
  p^*(c)=\frac{\varepsilon}{\varepsilon-1}c
\]
and the gross operating present value
\begin{equation}
  R(q,c)
  =
  \frac{Aq}{1-\beta\varphi}
  \frac{(\varepsilon-1)^{\varepsilon-1}}
       {\varepsilon^\varepsilon}
  c^{1-\varepsilon},
  \qquad
  \beta\in(0,1),\quad \varphi\in[0,1).
  \label{eq:p16m-return}
\end{equation}
The appendix derives the FOC, SOC, and accounting. In particular,
\(R_q>0\) and \(R_c<0\).

Internal production uses marginal cost \(c_I(m,k_i)\), with
\(c_{I,m}>0\) and \(c_{I,k}<0\), and setup cost \(F_I(m,k_i)\), with
\(F_{I,m}>0\) and \(F_{I,k}<0\). Both are finite for every \(k_i>0\). A
developer with little pre-existing manufacturing capability is not assumed to
be permanently incapable of internal production. It can in principle build,
acquire, validate, or adapt compliant capacity, but doing so is costly.
Accordingly, \(F_I(m,k_i)\) includes project-specific building, acquisition,
validation, adaptation, production-readiness, quality-system, and
organizational costs; it may be very large at low \(k_i\). Its route value is
\begin{equation}
  W_i^I(q,m)
  =
  s(q)R\!\left(q,c_I(m,k_i)\right)-F_I(m,k_i).
  \label{eq:p16m-internal}
\end{equation}

Entrusted production uses external marginal cost \(c_E(m)\), requires
\(b(m)>0\) units of qualified capacity with \(b'(m)>0\), incurs setup cost
\(F_E(m)\), and leaves the holder with burden \(\mu_E\geq0\). At the
qualified-capacity price \(p_m\),
\begin{equation}
  \begin{aligned}
  W_i^E(q,m;M,p_m)
  ={}&
  s(q)R\!\left(q,c_E(m)\right)-F_E(m)-p_mb(m)\\
  &-\mu_E-\tau_E(M).
  \end{aligned}
  \label{eq:p16m-entrusted}
\end{equation}
Entrusted production avoids establishing full proprietary manufacturing
capability but still requires technology transfer, validation, contracted
qualified capacity, and holder-side quality oversight. Each real monetary
cost enters once.
The transfer and abandonment values are
\begin{equation}
  W^T(q,m)=T(q,m),
  \qquad
  W^A=0.
  \label{eq:p16m-outside}
\end{equation}

The baseline therefore compares real organizational costs and manufacturing
capabilities without a route-financeability screen. Pre-MAH unavailability of
\(E\) arises only from \(\tau_E(0)=+\infty\). The appendix studies imperfect
finance as an inactive amplification extension.

\subsection{Organization and project advancement}
\label{subsec:p16m-organization}

At a fixed \(p_m\), the project chooses deterministically:
\begin{equation}
  \begin{aligned}
  W_i(q,m;M,p_m)
  &=
  \max\{W_i^I,W_i^E,W^T,W^A\},\\
  r_i^*(q,m;M,p_m)
  &\in
  \arg\max_{r\in\{I,E,T,A\}} W_i^r.
  \end{aligned}
  \label{eq:p16m-route-choice}
\end{equation}
Continuous heterogeneity makes ties a zero-measure event; no logit or
inclusive value is used.

The internal-minus-entrusted gap is
\begin{equation}
  \begin{aligned}
  \Delta_{IE}(k_i)
  ={}&
  s(q)\!\left[
  R\!\left(q,c_I(m,k_i)\right)-R\!\left(q,c_E(m)\right)
  \right]-F_I(m,k_i)\\
  &+F_E(m)+p_mb(m)+\mu_E+\tau_E(M).
  \end{aligned}
  \label{eq:p16m-gap}
\end{equation}
Holding \((q,m,M,p_m)\) fixed,
\[
  \Delta_{IE,k}
  =
  s(q)R_c\!\left(q,c_I\right)c_{I,k}-F_{I,k}>0.
\]
\begin{proposition}[Organizational sorting]
\label{prop:main-organizational-sorting}
Suppose \(\Delta_{IE}\) is continuous and crosses zero between two capability
levels. Because \(\Delta_{IE,k}>0\), there is a unique conditional cutoff
\(k^*(q,m;p_m,M)\) satisfying \(\Delta_{IE}(k^*)=0\). Developers below it
select \(E\), and those above it select \(I\), conditional on both retained
routes dominating \(T\) and \(A\). If transfer or abandonment dominates, the
cutoff does not determine the realized route.
\end{proposition}

Developers with weaker existing manufacturing capability face higher internal
variable and setup costs. Retained entrusted production is therefore
relatively more valuable for them, although transfer or abandonment may still
dominate both retained routes.

Define the best pre-MAH value and post-MAH value by
\begin{equation}
  \begin{aligned}
  W_i^0(q,m)
  &=\max\{W_i^I(q,m),W^T(q,m),0\},\\
  W_i^1(q,m;p_m)
  &=\max\{W_i^0(q,m),W_i^E(q,m;1,p_m)\}.
  \end{aligned}
  \label{eq:p16m-old-new-values}
\end{equation}
The MAH-relevant set is
\begin{equation}
  \mathcal C_i(p_m)
  =\left\{(q,m):
  W_i^E(q,m;1,p_m)>W_i^0(q,m)
  \right\}.
  \label{eq:p16m-relevant-set}
\end{equation}

\begin{proposition}[MAH-relevant project set]
\label{prop:main-relevant-set}
At a fixed \(p_m\),
\begin{equation}
  W_i^1-W_i^0
  =
  [W_i^E-W_i^0]_+,
  \label{eq:p16m-positive-part}
\end{equation}
so a strict project-value gain occurs exactly on
\(\mathcal C_i(p_m)\). Outside this set, the newly available route does not
improve on the best old option and the direct project-value effect is zero.
\end{proposition}

Expected optimized value per planning-stage project is
\begin{equation}
  \Omega_i(M,p_m)
  =
  \int W_i(q,m;M,p_m)\,dF(q,m).
  \label{eq:p16m-omega}
\end{equation}
Developer \(i\) solves
\begin{equation}
  \max_{x_i\geq0}
  \left\{
  \beta a_ix_i\Omega_i(M,p_m)
  -
  \frac{\kappa}{1+\nu}x_i^{1+\nu}
  \right\},
  \qquad \kappa>0,\quad \nu>0.
  \label{eq:p16m-advancement-problem}
\end{equation}
Strict concavity gives the unique solution
\begin{equation}
  x_i^*(M,p_m)
  =
  \left[
  \frac{\beta a_i}{\kappa}\Omega_i(M,p_m)
  \right]^{1/\nu},
  \label{eq:p16m-advancement}
\end{equation}
including \(x_i^*=0\) when \(\Omega_i=0\). The binary reform affects
\(x_i^*\) only through \(\Omega_i\). At a fixed \(p_m\), its effect is
strict only if \(E\) improves on the pre-MAH maximum on a
positive-probability project set.

\begin{proposition}[Project advancement and heterogeneity]
\label{prop:main-advancement}
At a fixed \(p_m\), the reform strictly raises \(x_i^*\) if and only if
\begin{equation}
  \int_{\mathcal C_i(p_m)}
  [W_i^E(q,m;1,p_m)-W_i^0(q,m)]\,dF(q,m)>0.
  \label{eq:p16m-strict-advancement}
\end{equation}
Advancement capability \(a_i\) scales any positive level response. Lower
\(k_i\) raises internal variable and setup costs, making entrusted production
relatively more valuable, but weak capability alone does not guarantee a
treatment effect because \(T\), \(A\), or a low entrusted value may still
dominate.
\end{proposition}

\subsection{Qualified manufacturing-service equilibrium}
\label{subsec:p16m-cmo}

Qualified supplier \(j\), with efficiency \(z_j\), chooses capacity
\(s_j\geq0\):
\begin{equation}
  \max_{s_j\geq0}\{p_ms_j-\Psi(s_j;z_j)\}.
  \label{eq:p16m-supplier}
\end{equation}
With \(\Psi_{ss}>0\), positive-price capacity is uniquely determined by
\(p_m=\Psi_s(s_j^*;z_j)\), and aggregate supply is
\[
  S_m(p_m)=\int s_j^*(p_m,z)\,dH_C(z).
\]

Expected entrusted capacity per planning-stage project is
\begin{equation}
  \chi_i^E(p_m;M)
  =
  \int b(m)\mathbf 1\{r_i^*(q,m;M,p_m)=E\}\,dF(q,m).
  \label{eq:p16m-project-demand}
\end{equation}
Study-related and total demand are
\begin{equation}
  \begin{aligned}
  D_m^{\mathrm{MAH}}(p_m;M)
  &=
  \int a_ix_i^*(M,p_m)\chi_i^E(p_m;M)\,dH(a,k),\\
  D_m(p_m;M)
  &=D_m^B(p_m)+D_m^{\mathrm{MAH}}(p_m;M).
  \end{aligned}
  \label{eq:p16m-demand}
\end{equation}
The demand expression includes both feedbacks: price changes expected value
and hence advancement, and it changes deterministic route selection.
Continuous heterogeneity makes aggregate demand continuous even though
individual indicators can jump.

The qualified-capacity market clears at the unique price
\begin{equation}
  D_m(p_m^*;M)=S_m(p_m^*).
  \label{eq:p16m-clearing}
\end{equation}
Supply is strictly increasing and demand weakly decreasing. Positive
background demand gives excess demand at zero, while entrusted and
background demand vanish and supply diverges at high prices. These conditions
give existence and uniqueness. If post-MAH study demand is positive at the
old price, \(p_m^*(1)>p_m^*(0)\): CMO scarcity attenuates the reform rather
than amplifying it through a direct price reduction.

For the equilibrium comparison, let \(p_m^h\equiv p_m^*(h)\), and define
\begin{align}
  \Delta\Omega_i^{\mathrm{dir}}
  &\equiv \Omega_i(1,p_m^0)-\Omega_i(0,p_m^0),
  &
  \Delta\Omega_i^{\mathrm{eq}}
  &\equiv \Omega_i(1,p_m^1)-\Omega_i(0,p_m^0),
  \label{eq:p16m-value-comparisons}\\
  \Delta x_i^{\mathrm{dir}}
  &\equiv x_i^*(1,p_m^0)-x_i^*(0,p_m^0),
  &
  \Delta x_i^{\mathrm{eq}}
  &\equiv x_i^*(1,p_m^1)-x_i^*(0,p_m^0).
  \label{eq:p16m-advancement-comparisons}
\end{align}

\begin{proposition}[CMO equilibrium and scarcity attenuation]
\label{prop:main-cmo}
Under atomless project and developer heterogeneity, the existing integrability
conditions, weakly decreasing demand, and strictly increasing supply, the
scalar CMO market has a unique clearing price. Moreover,
\begin{equation}
  p_m^*(1)\geq p_m^*(0),
  \qquad
  0\leq\Delta\Omega_i^{\mathrm{eq}}
  \leq\Delta\Omega_i^{\mathrm{dir}},
  \qquad
  0\leq\Delta x_i^{\mathrm{eq}}
  \leq\Delta x_i^{\mathrm{dir}}.
  \label{eq:p16m-scarcity-attenuation}
\end{equation}
The price inequality is strict when post-MAH study demand is positive at the
old price. MAH changes neither CMO supply nor the background demand schedule.
\end{proposition}

The baseline partial equilibrium contains exactly
\begin{equation}
  \left\{
  p_m^*,\;
  x_i^*,\;
  r_i^*(q,m),\;
  s_j^*
  \right\}_{i,j}.
  \label{eq:p16m-equilibrium}
\end{equation}
No entry, labor, capital, product-market aggregate, or welfare closure is
added.

\subsection{Observed outcomes and model boundary}
\label{subsec:p16m-predictions}

\paragraph{Planning, organization, and observed outcomes.}
At equilibrium,
\[
  \Lambda_i^{\mathrm{plan}}=a_ix_i^*,
\]
while expected realized retained outcomes are
\begin{equation}
  Y_i^{\mathrm{ret}}
  =
  a_ix_i^*
  \int s(q)
  \left[
  \mathbf 1\{r_i^*=I\}
  +
  \mathbf 1\{r_i^*=E\}
  \right]dF(q,m).
  \label{eq:p16m-retained-outcome}
\end{equation}
Thus observed outcomes can change through advancement or route composition.
Holder--producer separation is observed only after route assignment and
cannot cause the earlier \(x_i\) choice.

\paragraph{Novelty composition.}
For \(g\in\{O,\mathrm{Inc}\}\), outcomes use the same \(x_i^*\) and the
mixture \(\rho_gF_g\). The baseline permits either class to have a zero or
larger response and imposes no original-versus-incremental ranking. Patent
applications are not a baseline endogenous outcome.

\paragraph{Financing extension boundary.}
The baseline does not model financial intermediation or impose a binding
liquidity screen. Financing frictions may amplify the value of replacing a
large internalization commitment with a contract-based route; the technical
appendix analyzes that amplification separately and shows how the extension
nests this baseline when liquidity never binds.

\section{Empirical Predictions and Mapping Boundaries}

The deterministic model yields three conditional predictions. First, among
projects for which internal and entrusted production dominate transfer and
abandonment, developers with weaker internal manufacturing capability are
more likely to select retained entrusted production. Second, MAH raises
project advancement only when the entrusted route exceeds the best pre-reform
option; the response is zero when internalization, transfer, or delay remains
weakly better. Third, reform-induced demand for qualified manufacturing
services can raise \(p_m^*\), attenuating fixed-price gains in expected
project value and advancement. These predictions concern commercialization
organization and project advancement, not patent generation or scientific
productivity.

The empirical distinction is between ex ante advancement and downstream
observations. Product-level holder--manufacturer separation can proxy route
\(E\) only when holder identity, manufacturer identity, product identity, and
effective date are aligned. Approval or launch records are realized-product
outcomes, not direct observations of \(x_i\). Clinical milestones may proxy
project advancement only when their status and timing are measured consistently.
Patent applications are a separate patenting margin and are not a baseline
endogenous outcome. They may be studied as a separate empirical outcome, but
the present mechanism neither maps them to \(x_i\) nor assigns their reform
response a sign.

\begin{table}[htbp]
\centering
\footnotesize
\caption{Model Objects and Measurable Empirical Interfaces}
\label{tab:model_data_mapping}
\begin{tabularx}{\textwidth}{L{0.24\textwidth}Y Y}
\toprule
Model object & Candidate interface & Claim boundary \\
\midrule
Predetermined capability \(a_i\) & Pre-policy development history or other validated capability measure & Does not separately identify \(a_i\), and no policy shift in \(a_i\) is allowed. \\
Advancement \(x_i\) and \(a_ix_i\) & Consistently measured development-stage activity and planning-stage project events & \(x_i\) is not pure clinical effort, patent applications, or an approval count. \\
Internal manufacturing capability \(k_i\) & Pre-policy manufacturing assets, production licences, manufacturing history, validated technical scope, or other capability measures & This is the primary baseline heterogeneity; licence roles alone do not identify project-specific capability or realized route choice. \\
Route \(r_i^*=E\) & Product-level holder differs from manufacturer, with time-consistent identifiers & The observed split is downstream evidence on retained organizational assignment, not the policy treatment itself. \\
Retained outcome \(Y_i^{\mathrm{ret}}\) & Approval or launch with developer/holder identity & Combines advancement, project mix, deterministic route assignment, and \(s(q)\); it does not identify these primitives separately. \\
CMO scarcity \(p_m^*\) & Service price, or consistently measured capacity, utilization, density, waiting-time, or participation proxy & A nonprice proxy does not identify the price or the supply and demand schedules separately. \\
Novelty class \(g\) & Versioned regulatory classification \(O\) or \(\mathrm{Inc}\) & The classifier creates no \(x_{ig}\) and imposes no response ranking. \\
\bottomrule
\end{tabularx}
\end{table}

The timing restriction can be summarized as
\begin{equation}
  M
  \longrightarrow
  \text{anticipated availability and value of }E
  \longrightarrow
  \Omega_i
  \longrightarrow
  x_i^*,
  \label{eq:empirical-anticipated-chain}
\end{equation}
followed by
\begin{equation}
  \begin{aligned}
  x_i^*
  &\longrightarrow
  \text{planning-stage projects}
  \longrightarrow r_i^* \\
  &\longrightarrow
  \text{observed holder--producer separation}
  \longrightarrow
  \text{realized products}.
  \end{aligned}
  \label{eq:empirical-realized-chain}
\end{equation}
Observed separation cannot be placed before or treated as a cause of the ex
ante \(x_i\) choice. The original and incremental classes use the same
\(x_i\) and the mixture \(\rho_gF_g\); either class may have a zero or larger
response.

These entries are interfaces for future validated data. No data set is
asserted here to be available, complete, or identifying. Before use, a source
must pass provenance, coverage, key-uniqueness, unit, timing, and
construct-validity checks. Entry, route-specific realization, and smooth
route choice remain undeveloped extensions.

\subsection{Extension-Based Heterogeneity: Financing Capacity}

Financing is a secondary, extension-based heterogeneity test rather than a
baseline identifying restriction. Candidate measures must predate reform
exposure and may include registered or paid-in capital, cash and liquid
assets, leverage or debt capacity, listed status, VC/PE backing,
parent-company backing, or prior external financing. The empirical sequence
is to establish first whether weak internal manufacturing capability predicts
post-MAH retained entrusted production, then whether that route predicts
project advancement or retained realization, and only afterward whether
pre-policy financing remains informative after controlling for capability,
size, age, ownership, group affiliation, listed status, and innovation
history. A triple interaction among MAH exposure, weak manufacturing
capability, and financing capacity is therefore conditional and secondary,
not the model's primary baseline prediction. \citet{HallLerner2010} motivates
the plausibility of this extension; MAH is not modeled as a credit-supply
shock.

\section{Conclusion}

MAH changes expected project value by making retained entrusted manufacturing
available while preserving authorization-holder responsibility. The route is
most relevant when a developer can advance valuable projects but lacks
project-compatible internal manufacturing capability. Weak capability raises
the real variable and setup costs of internalization; the new retained route
can replace that large, lumpy requirement with contracted qualified
production. Its equilibrium value also depends on qualified-service scarcity:
reform-induced entrusted demand can raise \(p_m^*\) and attenuate the
fixed-price gain. Project advancement rises only for the MAH-relevant set and
need not rise for every developer or novelty class. Observed holder--producer
separation and realized products occur after advancement and organizational
assignment. Financing frictions can strengthen the mechanism when large
internalization commitments are difficult to fund, but the baseline does not
require a binding financing constraint. Patent applications and upstream
scientific research remain outside the baseline, and the model makes no
welfare claim.

\clearpage
\bibliographystyle{aer}
\bibliography{mah_route_indicator_friction_refs}

\end{document}
